\documentclass[11pt,a4paper,fontset=fandol]{ctexart}
\usepackage[margin=22mm]{geometry}
\usepackage{amsmath,amssymb,booktabs,longtable,array,tabularx}
\usepackage{enumitem,setspace,xcolor,hyperref,fancyhdr,needspace}
\hypersetup{colorlinks=true,linkcolor=blue!45!black,citecolor=blue!45!black,urlcolor=blue!45!black,pdftitle={Table I 共享估算基线 v0.2}}
\setstretch{1.13}
\setlength{\emergencystretch}{3em}
\setlength{\headheight}{14pt}
\setlength{\tabcolsep}{4pt}
\renewcommand{\arraystretch}{1.22}
\setlist{nosep,leftmargin=2em}
\pagestyle{fancy}
\fancyhf{}
\fancyhead[L]{\small Table I 共享估算基线 v0.2}
\fancyhead[R]{\small 中文研究底稿 · 待审阅}
\fancyfoot[C]{\thepage}
\newcommand{\RI}{\mathrm{RI}}
\newcommand{\sid}[1]{\texttt{#1}}
\newcommand{\ns}{\,\mathrm{ns}}
\newcommand{\us}{\,\mu\mathrm{s}}
\newcommand{\GBs}{\,\mathrm{GB/s}}
\title{\bfseries Table I 共享估算基线 v0.2\\[4pt]\large 28 nm CMOS 外围参考与两路吞吐估算方法}
\author{}
\date{2026 年 9 月 10 日 · 待审阅}
\begin{document}
\maketitle
\vspace{-1.5em}
\noindent 本底稿建立后续介质分析共同使用的参考设计与计算入口。研究对象是局部 CIM 服务单元，依次讨论共同外围、完整服务及两路吞吐；本版本止于这两节。
\section{28 nm CMOS 外围参考条件与通用开销}

\subsection{估计对象与共同条件的层次}

我们要回答一个参考设计问题：将不同介质的局部阵列配以共同的 28 nm CMOS 外围，在声明逻辑运算、内部展开和更新模式后，能够提供多大的 streaming 输入服务能力与 resident 写入能力。证据可以来自互补的器件、外围及宏级论文，组装后的结果由本文推导。

服务边界包含一个逻辑矩阵、实现其权重所需的局部阵列组合，以及本地输入寄存器、驱动、读出、数字重构、输出寄存器和一个更新控制域。八个权重位平面共同表示一份矩阵；互补存储若必需，也归入这份矩阵的实现。独立服务单元数固定为一，不计跨 bank、plane、die 的复制加速。输入从本地宽向量接口被接收，规定输出在本地结果寄存器就绪时完成；上游传输和下游搬运位于此边界之外。resident 服务则从本地写端口接收事务开始，到相应逻辑状态可计算为止。

参数分为三层。\textbf{A 类共同条件}规定逻辑尺寸、格式、边界、CMOS 数字节拍与局部写入资源；\textbf{B 类路径条件}分别规定 ACIM 的激励与 ADC 资源、DCIM 的数字计算组织；\textbf{C 类介质输入}提供物理响应、专用驱动、写入状态转换、读后恢复及刷新。共同 28 nm 是外围参考选择，原文工艺、电压和操作终点保持原样，器件物理时间也保留其自身条件。

\subsection{v0.1 推荐配置与输出约定}

推荐采用 $n_{\rm out}\times n_{\rm in}=128\times128$ 的逻辑矩阵，streaming 与 resident 均为二进制补码 INT8。沿用 Task 0 的单位，$b_S=b_R=1\,\mathrm{Byte/element}$；位数分别为 $8b_S=8b_R=8$。一个输入向量的 payload 是 128 Byte，整份权重的逻辑容量是 16384 Byte。容量只说明状态规模，一次更新的 $B_R$ 由更新事务实际完成的内容决定。

选择 128 行，是因为库内有直接的 128-cell 负载建模和 128 路输入宏；128 个输出使矩阵与逻辑字节换算简单。\sid{SACIM-03} 的物理 $128\times256$ 二进制阵列只对应 64 个四位权重输出；\sid{SDCIM-01} 的 $128\times128$ bit 阵列对应 16 个八位权重输出。本文的 $128\times128$ 明确指\emph{逻辑元素}，因此二进制映射需 $128\times128\times8=131072$ 个数据 bit，而不是直接照抄原宏的物理尺寸。\cite{CMOS-03,SACIM-03,SDCIM-01}

\begin{table}[htbp]
\centering\small
\begin{tabularx}{\textwidth}{@{}p{27mm} X@{}}
\toprule
项目 & 推荐选择与用途（均为本文参考选择）\\
\midrule
局部逻辑服务 & $128\times128$ INT8 矩阵乘一个 INT8 输入向量；交付 128 个 24-bit 结果寄存器。\\
CMOS 工作条件 & 28 nm，标称参考 0.9 V、25\,$^\circ$C；采用近标称工作模式的工程时隙范围。介质专用电压、电平转换另计。\\
输入与权重展开 & ACIM 每次输入 1 bit，128 路二电平驱动；8 个二进制权重位平面并行，平面内每列累加 128 项。\\
ACIM 读出共享 & 每个位平面 16 个 SAR，共 128 个；每个 ADC 在本平面复用 8 个输出列。每个共享组重新建立、求值和转换。\\
读出与重构 & 名义 10-bit SAR，约 8-bit 有效精度目标，归一到 8-bit 部分和码；16 条输出重构通道，每批用两拍进行位权合并和累加。\\
DCIM 组织 & 输入 1 bit/拍、权重 8 bit 并行；32 个输入项/组，16 个输出通道；输入位、行组、输出组顺序处理。\\
本地写入资源 & 一个更新域，128 路物理编码 bit 数据接口；每事务准备与完成各一拍。脉冲驱动的实际并行能力由介质填写。\\
主时序 & 各必要阶段顺序执行，同一单元不预设求值与写入重叠；向量捕获及最终输出提交合计两拍。\\
\bottomrule
\end{tabularx}
\caption{共同配置提案。ADC 数量与 DCIM 通道数是独立资源选择，未施加等面积约束。结构依据见正文，参数入口见共享 JSON。}
\label{tab:configuration}
\end{table}

这里将逻辑位宽与读出精度分开。对一对二进制输入、权重位平面，128 项部分和为 $0\ldots128$，需要 129 个可区分值，8-bit 码可容纳它们。完整 INT8 运算按位权重构，补码符号位使用负权 $-2^7$。24-bit 容器能容纳 128 项 INT8 乘积的整数和，并为中间重构保留空间。\textbf{DCIM 采用精确整数运算；ACIM 采用经校准的近似部分和及数字重构}。有效精度约 8 bit 是转换器目标，阵列误差和量化误差仍会传至最终结果；24-bit 输出容器只规定交付格式。本轮估计完成该服务的时间，算法精度验证留待需要精确性比较时补做。

\Needspace{5\baselineskip}
选择二电平输入和二进制权重位平面，是为了使公共主参考不依赖某种介质必须具备多级状态或高线性多位 DAC。八个位平面同时参与计算属于逻辑精度的资源成本。若介质需多级 cell、差分 cell 或更小子阵列，应提交等效映射与额外轮数；若另选近似模拟多级权重模式，应单列其精度及基线例外。

\subsection{原始证据如何形成外围范围}

\textbf{ADC 采用完整可复用时隙。} \sid{CMOS-02} 是直接的 28 nm 差分异步 SAR 实测依据，名义 10 bit，在 0.9 V、室温下以 100 MS/s 工作，测试结构含输入 buffer 与参考电路。对应的采样周期为
\[
T_{\rm sample}=\frac{1}{100\times10^6}=10\ns .
\]
这个 10 ns 是相邻采样的周期，覆盖采样与逐次比较过程，并非一次 SAR 比较器决策时间。原文没有给出适配当前 CIM 列负载后的独立转换定值。其 1 MHz 输入条件下正文给出 SNDR 55.13 dB、ENOB 8.86 bit；Fig.~10 的 SNDR 标作 56.91 dB，与正文不一致。奈奎斯特输入下报告 SNDR 51.54 dB、ENOB 8.27 bit，Table~2 同样列出后者。因此本参考使用\emph{约 8 bit 有效精度}的量级，不把名义 10 bit 当作十位无误精度。PVT 表来自后仿真，单独保留。\cite{CMOS-02}

参考时隙 $T_A$ 取 10--50 ns，推荐 20 ns。10 ns 保留独立 SAR 的采样周期锚点；20 ns 为移植到复用 CIM 读出的工程预算；50 ns 留出列选择、输入共模/参考适配、采样与局部锁存的宽裕预算。这些倍数是\emph{本文选择}，并非原文测得的 CIM 降速系数。用名义十位 SAR 得到约八位有效读出时，不再因为最终保留八位而按 $8/10$ 缩短周期。范围也不从各论文 headline 的最大、最小 rate 拼接而成。专用跨阻放大、积分或更慢的阵列建立属于 C 类输入。

每个 ADC 的时隙对应一个输出样本，一批 128 个 ADC 并行给出 128 个码，实际为 16 个逻辑输出各自的 8 个权重位平面部分和。一个输入 bit 需 8 个共享组，完整向量需 64 批。共享因子显式来自 $128/16=8$，而不是 ADC 面积与 cell-column 面积之比。\sid{CMOS-02} 的 0.026 mm$^2$ 活性区还包含 buffer/reference，直接复制 128 套该面积也不是当前布局评估；本轮选择的是转换资源预算，面积优化另议。

\textbf{输入采用数字二电平驱动。} 原文中的独立多位 DAC、数字字线激励和电荷耦合网络承担不同功能。\sid{SACIM-03} 使用 128 个四位输入 DAC，四位译码器选择 16 档电压，并经历 reset、evaluation、ADC readout；其宏在 100 MHz 下运行，提供约 10 ns 的\emph{整体}时序检查。p.~3 和 Table~II 采用 0.9 V，结论写作 1.2 V，本文保留这一差异，只取结构及 10 ns 周期量级。\sid{SACIM-05} 则将八位输入分成四组两位、经四电平输入并行求值，16 项累加的宏级 $T_{AC}$ 为 3.8 ns（0.9 V）或 3.6 ns（1.0 V）。其电压叠加有初始化、采样、叠加三阶段，读出后还进行数字移位累加。69 个 ADC 来自原宏的同位权合并，不是普遍的共享系数。\cite{SACIM-03,SACIM-05}

这些实测宏说明，在相应局部负载、低位输入和专门集成下，输入形成可以包含在数 ns 到十 ns 级宏服务内。本文选 $T_I=2$--10 ns、推荐 5 ns，覆盖输入码选择、分布式缓冲、公共开关复位及二电平激励就绪。2 ns 是简化为一位输入后的积极预算，10 ns 保留较宽的扇出与控制裕量；两端均为工程选择，\emph{没有从宏周期中减去假定 cell 延迟}。128 路输入在八个位平面间广播，使用各平面的本地缓冲；大阵列 RC、介质积分、特殊驱动电压的建立由后续介质增加。当前资料足以支持这一主输入方式，无需补一颗不匹配负载的独立八位 DAC。

\textbf{数字路径采用节拍加完整周期数。} \sid{SDCIM-03} 在 0.8 V 下给出 3 ns 计算时钟，八位输入分八拍，并有 32 项 post-sum。\sid{SDCIM-01} 报告 1.1 V 下约 360 MHz（Fig.~10 标 363 MHz），即约 2.8 ns/拍；其 $128\times16$ INT8 VMM 实际需 64 拍，包含输入位展开和空间分组。\sid{SDCIM-05} 的 INT4 一 MAC 拍、INT8 两 MAC 拍；Table~III 的 0.9 V accurate INT4/INT8 access time 为 1.0/2.1 ns，作为更并行结构的交叉依据。该文 FP8 的正文与表格互换了 2.5/2.7 ns 的 AP/AC 标签，故本文不据此选择 INT8 时钟。\cite{SDCIM-01,SDCIM-03,SDCIM-05}

共同数字节拍取 $T_D=2$--10 ns，推荐 5 ns（200 MHz）。2 ns 是具备相应并行与流水资源时的积极设计端；5 ns 围绕几个 ns 的硅测时序留出移植余量；10 ns 为较大局部归约、布线和控制负载保留空间。文献中的更低电压模式可明显慢于此范围，本参考固定近标称模式。来自 0.8、1.0、1.1 V 的锚点不被宣称为固定 0.9 V 的 PVT 保证，也不作电压或工艺线性缩放。

DCIM 一组运算的公共数字阶段包含输入位选择、局部乘法/归约、部分和更新和寄存控制；阵列读侧物理阶段另有定义。ACIM 一批的数字重构分两拍：第一拍合并八个位平面的正负位权，第二拍按输入位权更新 16 个输出累加器。这里的通道数与拍数是资源预算，后续实现若需要更深的归约流水，应显式增加拍数。数字节拍也用于公共命令与完成接口，统一传播其不确定性。

\textbf{写入只共享低压接口及控制预算。} \sid{CMOS-03} 的 28 nm SRAM 分析采用 TT、1 V、50 FO4 约 1 ns 周期、25 FO4 字线脉宽和 15 fF 的 128-cell 位线负载，并以读/写成功及连续操作终点定义周期；这些是模型条件。\sid{CMOS-07} 的 28 nm eFlash 平台 SRAM 有 48-bit 字、1024 字、列 MUX4；Table~II 给出典型 read access 1.54 ns 与 455 MHz 时钟，后者约 2.20 ns。两者共同支持数 ns 级局部普通读写/控制量级，同时显示 access、WL 脉冲与可重复周期是不同量。该 eFlash 工艺具有较高阈值，其数据不被推广为全部逻辑 compiler 的定值。\cite{CMOS-03,CMOS-07}

参考采用 128-bit 物理编码数据接口，结构锚点来自 \sid{SDCIM-01}，而 \sid{SDCIM-05} 的普通端口只有 32 bit；二者都没有独立闭合的写周期。本文选择每事务准备与完成各一拍，$T_W=2T_D=4$--20 ns，推荐 10 ns。第一拍可接收一个 128-bit 数据批，额外数据批各占 $T_D$；page 接口若命令与数据分开，全部数据拍另计。该预算覆盖命令、数据锁存、低压驱动接口和完成状态接收，物理编程脉冲、读回 sensing/ADC、verify 比较和必要高压产生按实际序列加入。若原文给出完整普通写周期并已覆盖这些功能，采用完整周期替代相应块，保持原来已包含的边界。

\Needspace{24\baselineskip}
\subsection{核心范围表与完整服务开销}

% Generated by scripts/check_shared.py --emit.
\begingroup\small
\begin{longtable}{@{}>{\raggedright\arraybackslash}p{23mm}>{\raggedright\arraybackslash}p{49mm}>{\raggedright\arraybackslash}p{24mm}>{\raggedright\arraybackslash}p{58mm}@{}}
\caption{共同外围能力范围。每项按自身粒度使用，完整服务的次数由结构推导。}\label{tab:periphery}\\
\toprule 模块 & 服务粒度 & 范围/推荐 & 依据与用途\\\midrule\endfirsthead
\toprule 模块 & 服务粒度 & 范围/推荐 & 依据与用途\\\midrule\endhead
输入形成与公共复位 $T_I$ & 每次求值的公共复位、128 路二电平输入与本地缓冲。 & 2--10 ns\newline 推荐 5 ns & \sid{SACIM-03}, \sid{SACIM-05}, \sid{CMOS-03}。\newline 更改多位输入或扇出负载时核对驱动与精度；专用建立另计。\\
采样、转换与局部读出 $T_A$ & 一颗 ADC 一个样本的完整可复用时隙；并行颗数属于结构参数。 & 10--50 ns\newline 推荐 20 ns & \sid{CMOS-02}, \sid{SACIM-03}, \sid{SACIM-05}。\newline 名义 10 bit、有效约 8 bit；含 mux/采样/转换/局部锁存，跨批保持另计。\\
数字服务节拍 $T_D$ & 一个局部数字阶段；拍数由运算、归约与通道资源推导。 & 2--10 ns\newline 推荐 5 ns & \sid{SDCIM-01}, \sid{SDCIM-03}, \sid{SDCIM-05}。\newline 参考 200 MHz；共同近标称预算，所选负载下需相应资源。\\
写入准备与完成控制 $T_W$ & 每事务准备和完成两拍；编码数据装入另按接口规则计。 & 4--20 ns\newline 推荐 10 ns & \sid{CMOS-03}, \sid{CMOS-07}, \sid{SDCIM-01}, \sid{SDCIM-05}。\newline 首数据批可与命令同拍；完整更新周期已含接口时替代。\\
\bottomrule\end{longtable}
\endgroup


表中的四块足以组织主参考。ACIM 一次物理步骤经历输入形成、阵列求值、转换、两拍数字重构及必要恢复；$t_{m,A}$ 汇集其中公共预算以外的阵列建立/积分、专用驱动和读后恢复时间，各段按实际操作顺序执行。每个 ADC 共享组重新求值，因而不依赖所有模拟结果长期保持。八个输入 bit 乘八个共享组，共 64 步，向量捕获和最终输出提交另计两拍。因此\emph{纯公共外围占用}为
\begin{equation}
T_{S,A,\rm peri}=64(T_I+T_A+2T_D)+2T_D .
\label{eq:acim_peri}
\end{equation}
固定配置下，短、参考、长时隙情景分别为 1028、2250、5140 ns；展示为约 1--6\,$\mu$s，参考约 2.3\,$\mu$s。它是完整逻辑服务所需的公共外围占用，介质时间仍须加入，未产生任何介质的最终 $\rho$。

DCIM 的 8 个输入 bit、4 个行组、8 个输出组共 256 个计算步骤。每步一拍公共数字阶段，边界两拍，所以纯公共外围占用为 $258T_D=516$--2580 ns，展示约 0.5--3\,$\mu$s，参考约 1.3\,$\mu$s。这些完整服务量级比某些已发表宏的单周期长，主要来自所选逻辑尺寸、位展开和 16 通道复用；原宏的更大并行、模拟位权合并或不同累加维度并未被暗中移植进来。

本配置让计数清楚，但它不是所有介质的唯一合理实现。最值得审阅的是 128 个 ADC 的资源预算、16 条 DCIM 输出通道，以及主参考每个共享组重新求值的组织。这些选择会系统改变服务次数；冻结后，各介质沿用同一基线或明确列出例外。

\clearpage
\section{\texorpdfstring{局部 CIM 服务单元的 $\rho$、$\tau$、$\RI^*$ 估算方法}{局部 CIM 服务单元的 ρ、τ、RI* 估算方法}}

\subsection{从逻辑服务字节与服务间隔出发}

一个完整 streaming 服务单元对应 $B_S$ Byte 输入，完成规定求值与输出；一个完整 resident 更新服务单元对应 $B_R$ Byte，更新后已建立可计算的逻辑状态。它们分别以 $\Delta_S$、$\Delta_R$ 秒的服务间隔完成，则
\begin{equation}
\boxed{\rho\approx\frac{B_S}{\Delta_S},\qquad
\tau\approx\frac{B_R}{\Delta_R},\qquad
\RI^*=\frac{\rho}{\tau}\approx
\frac{B_S}{B_R}\frac{\Delta_R}{\Delta_S}.}
\label{eq:service}
\end{equation}
$\rho,\tau$ 的单位均为 Byte/s，$\RI^*$ 为同量纲服务字节的比值，数值无量纲。$B_S,B_R$ 是硬件事务粒度；Task 0 的 $Q_S,Q_R$ 则是在 workload 窗口内累计的逻辑输入与逻辑写入量，包含实际重放和重载。内部位展开、ADC 批次和编程重试只增加相应服务占用。

本参考 $B_S=n_{\rm in}b_S=128$ Byte。若一次局部更新真正完成 16 个 INT8 权重，$B_R=16$ Byte，式~\eqref{eq:service} 的字节因子为 8；若把同样的连续更新聚合为整份 16384 Byte 状态，则 $\Delta_R$ 也应包含完成整份状态的全部时间。以均匀、无额外固定开销的更新为例，分子和分母同比聚合不改变 $\tau$。因此 ridge 一般不等于“写延迟/读延迟”。相同单元独立复制 $k$ 份、两路需求和可用能力均按相同口径汇总时，$\rho$ 与 $\tau$ 同乘 $k$，ridge 可抵消这一因子；ADC 复用、写入并行宽度、不同位串行次数及专用电压准备通常只影响其中一路。

对窗口内 $Q_R>0$，接回已审阅模型：
\[
\RI=Q_S/Q_R,\qquad P=Q_S/T\leq\min(\rho,\tau\RI).
\]
$Q_R=0$ 使用静态分支 $\RI=\infty$、$P\leq\rho$。用本节的估计能力代入，得到的是对应参考情景的估计 Roofline；区间表示配置和参数的不确定性，不是统计置信区间，也不是全部实际实现的严格包络。两路可分别估计，即使共享阵列而互斥，原上界仍成立；能否同时达到其独立能力由实际调度决定。

\subsection{Streaming：把完整精度落实为物理步骤}

\textbf{ACIM 主参考。} 将 INT8 输入按 8 个 bit 展开，8 个二进制权重位平面同时求值。对每个输入 bit，每平面 128 个输出需由 16 个 ADC 分 8 批读出。因此一次向量有 64 次物理激励/求值、64 批转换，标量 ADC 结果总数为
\[
8\ \text{输入 bit}\times128\ \text{输出}\times8\ \text{权重位平面}
=8192=64\times128.
\]
这 8192 个码完成同一次输入服务，仍只计 $B_S=128$ Byte。每批服务 16 个输出，每个输出的八个权重位平面结果按正负位权合并，再按当前输入位权累加；所有 128 个输出就绪才完成服务。

令 $t_{m,A}$ 是本次物理步骤中介质读侧的额外必要占用，包括读激励后阵列 RC/积分/响应、公共接口以外的专用驱动及必要读后恢复，内部顺序按介质组织。主参考各阶段顺序执行，服务间隔直接为
\begin{equation}
\boxed{\Delta_{S,A}=64(T_I+t_{m,A}+T_A+2T_D)+2T_D .}
\label{eq:acim}
\end{equation}
各时间用相同单位后求和，代入式~\eqref{eq:service} 前统一转成秒。若各步骤物理时间不同，用步骤求和替代 $64t_{m,A}$。多个互补读必须分别计入 $t_{m,A}$ 或步数；若能一次差分 sensing，则按那一次实际服务计。八个位平面的并行已经用在步数中，吞吐不再额外乘八。

当某介质只能同时激活较少行、轮询权重平面、或使用更少 ADC，需显式改变行组数、权重组数和转换批次。广义计数可以按“输入分片数 $\times$ 行组数 $\times$ 权重组数 $\times$ 每组输出批数”组织，但各因子必须对应实际独立循环；同一个结构约束不同时在 ADC 数量和额外折扣中计两次。多位模拟输入或多级权重模式则同时改变可分辨部分和、误差条件和数字重构，按另一声明配置重算。

\Needspace{6\baselineskip}
\textbf{DCIM 主参考。} 公共实现一次处理 32 个输入项和 16 个输出，输入逐 bit、权重八位并行。每个数字计算步骤把当前输入位与权重相乘、归约 32 项，再将带输入位权的贡献加到该组 24-bit 输出累加器。4 个行组覆盖 128 项，8 个输出组覆盖 128 个结果，所以
\[
N_D=\underbrace{8}_{\text{输入 bit}}
\underbrace{\left\lceil128/32\right\rceil}_{\text{行组}}
\underbrace{\left\lceil128/16\right\rceil}_{\text{输出组}}=256.
\]
令 $t_{m,D}$ 为每组数字求值前、公共数字阶段以外的物理读侧时间及必要恢复，则
\begin{equation}
\boxed{\Delta_{S,D}=256(t_{m,D}+T_D)+2T_D .}
\label{eq:dcim}
\end{equation}
输入组选择和部分和寄存已在 $T_D$ 中，输入捕获和最终提交在两拍边界开销中。主参考完整精度计数为 256 步，不把每一拍当作服务了 128 Byte 新输入。若实现的权重也按位串行，应按实际独立循环增加权重步骤；若采用 2-bit/4-bit 输入、整字并行乘法或更多输出通道，应重算步数和所需时钟，而不是保持原拍数再额外乘并行收益。

文献宏若已有完整计算周期 $T_{\rm cycle,full}$，并覆盖对应的物理响应、感测、逻辑和恢复，则用它替代式~\eqref{eq:dcim} 中相应的整体步骤。所引用的逻辑维度、输入位宽和输出粒度还要匹配当前分组；例如 \sid{SDCIM-01} 的 64 拍服务的是其 $128\times16$ VMM，不能直接给当前 $128\times128$ 矩阵也填 64 拍。类似地，ACIM 的整体宏周期若已包含输入、ADC 和重构，按其实际模式换算完整服务后使用，原有阶段不再追加。采用整体证据前，还须核对其资源、精度和时隙与所选共同情景兼容；不兼容的实测宏用于独立交叉检查或显式例外，不据此静默覆盖共同节拍。

\textbf{主调度与流水线对照。} 式~\eqref{eq:acim}--\eqref{eq:dcim} 的 sum 来自已选择的顺序组织，故无跨向量重叠时服务间隔等于完整串行服务时长。若局部缓冲允许输入/阵列阶段、ADC 阶段、数字阶段流水，令每个逻辑向量对共享资源 $r$ 的必要占用为 $D_r$，则稳态间隔满足 $\Delta_S\geq\max_rD_r$。只有资源互相独立、缓冲足够且依赖允许时，才能将该值作为可实现的间隔估计；同一资源上多个阶段的占用先相加。转换延迟不等于 ADC initiation interval，流水线报告优先使用后者。有限向量数还包括流水线填充、排空。v0.1 数值采用串行公式；流水化若作为后续敏感性对照，另存成对情景。

\subsection{Resident：按完整更新序列计算}

\textbf{A. 可直接更新型。} 在声明模式下，一组权重经一次完整更新周期后可计算，则
\begin{equation}
\tau\approx\frac{B_{\rm update}}{T_{\rm update,full}}.
\label{eq:direct}
\end{equation}
先由映射求完成的逻辑权重数 $L_{\rm done}$，再取 $B_{\rm update}=L_{\rm done}b_R$。若每逻辑权重需要 $c$ 个 cell（已含二进制位切片、互补与编码），一次并行 $p$ 个 cell 完成完整目标状态，且分组对齐，可取 $L_{\rm done}=\lfloor p/c\rfloor$。若这些 cell 只是一个权重的部分位，则先累计完成所有必要位再确认逻辑 payload；不要按已触及的权重数记作全部完成。

主参考无冗余二进制映射中，一个权重需 8 个 cell。128 路驱动若能同时完成 128 个相应 cell，便完成 16 Byte；若高压/电流预算只支持较少并行 cell，增加局部写步骤或减小每事务 $B_R$。数据传输、命令、编程脉冲、建立及恢复共同构成 $T_{\rm update,full}$。对于一次数据批、单次完整物理更新的简化模式，可写 $T_{\rm update,full}=2T_D+t_{\rm physical,full}$；原文的完整更新周期若已包括公共准备与完成，直接使用该周期。

\textbf{B. 分步编程与独立擦除型。} 将实际状态转换组织为预处理、所需编程步骤、读回校验、恢复及适用擦除。例如
\[
T_{\rm update,full}=T_{\rm front}
+\sum_j\bigl(T_{{\rm drive+program},j}+T_{{\rm verify},j}
+T_{{\rm recover},j}\bigr)+T_{\rm erase,amort}.
\]
这些是粗粒度的操作块；相互重叠时按资源服务能力组织。SET 和 RESET 时间不同，只说明状态转换有方向性；每次事务执行哪一步、是否先全域 RESET，再选定 SET，取决于更新策略。初次形成若为部署初始化只在相应窗口计入；每次都需要的预处理则进入每次更新。

对于单个局部编程域的页/块模式，持续整块循环重写且每次擦除后使用 $N_{p/e}$ 个有效页服务时，可以采用
\begin{equation}
\tau\approx\frac{B_{\rm page,logical}}
{T_{\rm front}+T_{\rm program,full}+T_{\rm erase}/N_{p/e}}.
\label{eq:page}
\end{equation}
$T_{\rm front}$ 是原文 program cycle 外部仍需的命令、数据装入及完成接口时间。$B_{\rm page,logical}$ 是本逻辑映射在此页事务中完成的\emph{有效 resident Byte}；物理页未使用的字节和为旧数据保留的内容进入实际成本，不扩充逻辑写入分子。$N_{p/e}$ 指每次擦除真正支撑的页服务数，若只更新部分页，按其情景重新摊销。该式中的页容量不赋予整 die 的多 plane 并行；如果页天然覆盖当前矩阵以外的内容，本单元承担的页操作及有效字节必须同时说明。保持原生 page/block 粒度即可，无需设想器件支持逐字编程。

在已有预擦除空间上 append，可将该窗口的擦除项设为零，并明确预擦除初始状态；持续循环重写则计入擦除摊销。两种模式分别报告，避免混用 append 的时间与循环重写的容量供给条件。若映射跨多个物理页才能完成一组逻辑权重，用这些页的完整服务总时间与对应逻辑字节配对。

\textbf{C. 闭环校验。} 若原文报告从发起编程至目标状态/误差终点完成的 program cycle，且已含脉冲、verify/retry，直接以该完整周期作为操作块。若只报告单脉冲 $t_p$，则需写出脉冲数与每次读回/比较时间；例如固定 $K$ 次闭环且每次后都校验的简化情景为 $K(t_p+t_v+t_r)$，另加一次性准备。$K$ 由原文分布、终止规则或明确情景给出，末次是否需要校验按流程处理。分子仍为只完成一次的 $B_R$。读回用 binary sensing 还是多位 ADC、是否共享 streaming ADC、并行比较宽度多少，均应和同一局部资源配置衔接。

\subsection{维护、逻辑精度与局部边界的放置规则}

refresh 和破坏性读后的 restore 属于内部维护：与每次求值绑定的 restore 放入该步的物理服务时间；周期刷新占用同一资源时，可将已知刷新占用比 $f$ 作用为该资源可用率 $1-f$，或直接在排程中加入刷新块。两种方法选其一。若读写受到不同资源占用，则分别作用于两路；刷新数据不作为新 workload resident payload。endurance 和寿命单列为部署约束，不混入瞬时 $\tau$。

逻辑 INT8 不规定一个 cell 存八位。每个 cell 的状态位数、互补数量、串行步骤和可同时完成的逻辑权重数，一起决定 $B_R$ 与 $\Delta_R$。位切片并行已经进入逻辑并行度时，不再重复扩充 payload；重复物理步骤已经进入时间时，不再对 Byte 额外折扣。页/块容量和 3D 层数描述组织与容量，读写并行数由选通、驱动、感测及编程资源决定。

两路能力必须引用同一份局部状态实现、器件模式和外围资源预算。可以分别测算单独提供 streaming 或更新时的服务能力；同时运行争用同一 ADC、写驱动或阵列时，记录互斥/共享条件。若更新过程需要不同电压、重新校准或较长恢复，这些代价放在相应事务末端，使“完成更新”确实意味着能够继续计算。

\subsection{情景范围、两个示意算例与后续接口}

范围生成遵循固定次序：先固定逻辑配置、映射与服务模式；给共同外围及介质时间各自定义范围；枚举物理兼容的组合；在每个组合内导出 $(\rho,\tau,\RI^*)$；保存未取整值，再取适合量级分析的展示范围，并指出主要驱动参数。共同外围采用短时隙 $(T_I,T_A,T_D)=(2,10,2)$ ns、参考 $(5,20,5)$ ns、长时隙 $(10,50,10)$ ns 三个共享情景。它们是工程预算，未被命名为已验证的 PVT 角。各介质采用相同规则，物理时间变化不必与 CMOS 时隙同向；若存在独立变化且模式兼容，应加入相应组合。

ridge 优先取各成对情景内比值的最小、最大值；仅用独立吞吐区间交叉除法得到的
\[
\left[\rho_{\min}/\tau_{\max},\ \rho_{\max}/\tau_{\min}\right]
\]
标为保守外包络，区别于实际计算的成对范围。仅当共因子以相同乘法形式作用于两路时抵消；例如 $\Delta_S=C_S T_D$、$\Delta_R=C_R T_D$ 可消去 $T_D$，而 $(aT_D+t_{m,S})/(bT_D+t_{m,R})$ 通常保留其影响。

以下\textbf{两个算例均为不绑定介质的算术示意}，其物理时间为人为选定的输入，仅检验服务粒度、单位与范围传播。三个配对场景是演示路径，不是允许参数空间的穷尽。

\textbf{示意一：ACIM 与一次直接更新。} 采用式~\eqref{eq:acim}，三个情景的 $t_{m,A}$ 依次为 10、25、40 ns；一次完整物理更新为 20、50、100 ns，128 个二进制 cell 同时完成 16 个 INT8 权重，即 $B_R=16$ Byte。总更新时间为物理更新加 $2T_D$。参考情景得到
\[
\Delta_S=64(5+25+20+2\times5)+2\times5=3850\ns,
\quad \Delta_R=50+2\times5=60\ns,
\]
\[
\rho\approx0.0332\GBs,\quad\tau\approx0.267\GBs,\quad
\RI^*=\frac{128}{16}\frac{60}{3850}\approx0.125.
\]
这里省略字节因子会使 ridge 错八倍。ADC 复用和八位输入展开已在 64 步中，payload 不再乘 64。

\textbf{示意二：DCIM 与页更新。} 采用式~\eqref{eq:dcim}，三个 $t_{m,D}$ 为 3、10、30 ns。有效逻辑页为 256 Byte，映射恰好占 2048 物理 bit；128-bit 接口需 16 个数据拍，命令/完成另占两拍。设完整编程周期依次为 5、10、20\,$\mu$s，擦除周期为 100、200、400\,$\mu$s，每擦除支撑 16 页服务。此处完整编程含其内部 verify，时间不再乘校验次数。参考情景为
\[
\Delta_S=256(10+5)+2\times5=3850\ns,
\]
\[
\Delta_R=(16+2)\times5\ns+10\us+200\us/16
=22.59\us.
\]
故 $\rho\approx0.0332\GBs$、$\tau\approx0.0113\GBs$，ridge 约 2.93；服务字节因子在此为 $128/256=1/2$。页字节、页数和擦除时间的单位分别参与一次换算。

% Generated by scripts/check_shared.py --emit; illustrative, no medium binding.
\begin{table}[htbp]\centering\small
\begin{tabular}{@{}llrrrrr@{}}\toprule
示意 & 情景 & $\Delta_S$ (ns) & $\Delta_R$ (ns) & $\rho$ (GB/s) & $\tau$ (GB/s) & $\RI^*$\\\midrule
1 & 短时隙 & 1668 & 24 & 0.07674 & 0.6667 & 0.1151\\
1 & 参考 & 3850 & 60 & 0.03325 & 0.2667 & 0.1247\\
1 & 长时隙 & 7700 & 120 & 0.01662 & 0.1333 & 0.1247\\
\midrule
2 & 短时隙 & 1284 & 11286 & 0.09969 & 0.02268 & 4.395\\
2 & 参考 & 3850 & 22590 & 0.03325 & 0.01133 & 2.934\\
2 & 长时隙 & 10260 & 45180 & 0.01248 & 0.005666 & 2.202\\
\bottomrule\end{tabular}
\caption{两个算例的成对情景。时间列保留算术结果用于复算，非器件测量精度；GB 为 $10^9$ Byte。}
\label{tab:examples}\end{table}
示意 1 的展示范围为 $\rho$=0.016--0.077 GB/s、$\tau$=0.13--0.67 GB/s，成对 ridge 为 0.11--0.13；独立端点外包络为 0.024--0.58。
示意 2 的展示范围为 $\rho$=0.012--0.1 GB/s、$\tau$=0.0056--0.023 GB/s，成对 ridge 为 2.2--4.4；独立端点外包络为 0.55--18。


保留数据中的所有计算位数便于复算；示意范围按两位有效数字向外取整，正文外围整向量时间按一位有效数字向外取整。若只有少数情景，表内原计算点与取整范围同时呈现，避免用显示精度暗示物理准确性。示意一主要由 ADC 批次和写周期影响，示意二由数字步骤及 program/erase 摊销影响；改变普通外围的同一参数后，两路和 ridge 必须一起重算。

后续每类介质只需提交以下最小输入：所选器件/状态与逻辑映射；读侧响应及终点条件；直接更新或预处理、编程、擦除、校验、恢复的完整步骤与时间；实际局部读写粒度和并行约束；page/block 的有效映射；专用驱动和维护；引用的共享版本与例外；兼容的成对情景、导出范围及主导因素。来源记录保留原值、单位、模式、PDF 页码和图表。共享参数由 \texttt{data/shared\_parameters.json} 的 v0.1 入口提供，审阅后冻结；例外单列，避免覆盖共同条件。至此接口已具备，本轮不建立各介质结果章节。

\clearpage
\bibliographystyle{unsrt}
\bibliography{references}
\end{document}
